\documentclass[10pt,twocolumn,letterpaper]{article}
\usepackage{microtype}
\usepackage{booktabs}
\usepackage{hyperref}\usepackage[spanish]{babel}
\usepackage{graphicx}
\usepackage{amsmath,amssymb}
\usepackage{cleveref}
\usepackage{textcomp}
\usepackage{xcolor}
\raggedbottom

\title{Rigor Estadístico en Evaluación de Modelos YOLO: Intervalos de Confianza Bootstrap, Análisis de Modos de Fallo y Pruebas Wilcoxon}
\author{William Steve Rodriguez Villamizar (wisrovi rodriguez)\\AI Leader \& Solutions Architect\\wisrovi-suit (https://github.com/wisrovi/w-cli)}
\date{}

\begin{document}
\maketitle

\section{Abstract \& Keywords}
\textbf{Abstract:} Reporting a single mAP value after YOLO training conveys no information about whether observed differences between models are real or artifacts of random seed variation. We present a statistical validation pipeline that combines bootstrap resampling (1000 iterations, 95\% confidence intervals), failure mode analysis via FiftyOne, and Wilcoxon signed-rank tests to determine whether model improvements are statistically significant. Evaluated on five YOLO datasets comparing YOLOv8n against YOLO26n, our pipeline reveals that 3 of 5 apparent improvements (mAP gains of 0.3--1.1\%) fall within the bootstrap confidence interval and are not statistically significant ($p > 0.05$). Only 2 of 5 improvements survive the Wilcoxon test at $\alpha = 0.05$. The failure mode analyzer identifies high-confidence false positives (confidence $> 0.75$) as the dominant error category, accounting for 62\% of critical failures. An ablation study shows that omitting the bootstrap step leads to overconfident claims in 60\% of comparisons. The pipeline executes in 12 seconds on a single GPU and integrates into the existing post-training workflow as steps 11 and 13.

\textbf{Keywords:} Bootstrap Resampling, Statistical Significance, YOLO Evaluation, Failure Mode Analysis, Wilcoxon Signed-Rank Test, Confidence Intervals, MLOps.

\section{Author Information}
This research was conceptualized and developed by William Steve Rodriguez Villamizar (wisrovi rodriguez), AI Leader \& Solutions Architect for the wisrovi-suit ecosystem (https://github.com/wisrovi/w-cli). Contact: wisrovi.rodriguez@gmail.com.

\section{Introduction}
A YOLOv8n model achieves 91.3\% mAP$_{50}$ on a defect detection dataset. A YOLO26n model achieves 92.1\% on the same split. The paper concludes: ``YOLO26n outperforms YOLOv8n by 0.8\%.'' But is this difference real?

With a single random seed, we cannot answer this question. The 0.8\% gap could reflect a genuine architectural advantage---or it could fall within the noise floor of a single train/test split. Dem{\v{s}}ar~\cite{demsar2006statistical} demonstrated this problem comprehensively: comparing classifiers on single data splits without statistical testing is ``the most common methodological error in machine learning research.''

We address this with a three-component statistical validation pipeline:

\begin{enumerate}
    \item \textbf{Bootstrap Evaluator} (step 11): Computes 95\% confidence intervals via 1000 bootstrap resamples of the prediction set. If the CI for the accuracy difference includes zero, the improvement is not significant.
    \item \textbf{Failure Mode Analyzer} (step 13): Uses FiftyOne to load the test dataset, evaluate detections against ground truth, and extract high-confidence false positives (confidence $> 0.75$) and all false negatives. This enables targeted debugging.
    \item \textbf{Wilcoxon Signed-Rank Test}: A paired, non-parametric test applied across 5 cross-validation folds to determine if the distribution of per-fold accuracy differences is significantly different from zero.
\end{enumerate}

These components run sequentially in the post-training pipeline, consuming 12 seconds total. No manual statistical analysis is required.

\section{Related Work}
Statistical comparison of classifiers has been formalized by Dem{\v{s}}ar~\cite{demsar2006statistical}, who established the Friedman test followed by Nemenyi post-hoc analysis as the standard for multi-classifier comparison. For paired comparisons on a single dataset, the Wilcoxon signed-rank test~\cite{wilcoxon1945individual} is the non-parametric equivalent of the paired t-test, requiring no normality assumption.

Bootstrap confidence intervals were introduced by Efron and Tibshirani~\cite{efron1993introduction}. Their key insight: by resampling the observed data with replacement $B$ times, we can estimate the sampling distribution of any statistic without distributional assumptions. For classification accuracy, $B = 1000$ typically yields stable estimates~\cite{davison1997bootstrap}.

Failure mode analysis for object detection was formalized by Hoiem et al.~\cite{hoiem2012diagnostic}, who categorized errors into localization failures, confusion with background, class confusion, and missed detections. FiftyOne~\cite{fiftyone2024} operationalizes this framework for modern workflows.

The YOLO ecosystem~\cite{jocher2023yolo} typically reports single-seed metrics. Our pipeline addresses this gap by embedding statistical validation directly into the post-training workflow.

\section{Proposed Architecture / Methodology}

\subsection{Bootstrap Confidence Interval}
Given $n$ predictions $\hat{y}_1, \ldots, \hat{n}$ and targets $y_1, \ldots, y_n$, we compute bootstrap accuracy:

\begin{equation}
    \hat{\theta}^*_b = \frac{1}{n} \sum_{i=1}^{n} \mathbb{1}[\hat{y}^*_{b,i} = y^*_{b,i}], \quad b = 1, \ldots, B
\end{equation}

where $(\hat{y}^*_b, y^*_b)$ are drawn with replacement from the original set. The 95\% CI is:

\begin{equation}
    \text{CI}_{95\%} = \left[\hat{\theta}^*_{(0.025 \cdot B)}, \; \hat{\theta}^*_{(0.975 \cdot B)}\right]
\end{equation}

where $\hat{\theta}^*_{(p)}$ denotes the $p$-th percentile of the bootstrap distribution.

\subsection{Failure Mode Classification}
For each detection, we classify:

\begin{itemize}
    \item \textbf{False Positive (FP)}: predicted class $\neq$ true class, confidence $> 0.5$.
    \item \textbf{False Negative (FN)}: ground truth object not detected.
    \item \textbf{Critical FP}: FP with confidence $> 0.75$---these are the most dangerous errors.
\end{itemize}

Critical FPs are exported as a FiftyOne dataset for visual inspection.

\subsection{Wilcoxon Signed-Rank Test}
For $k$ cross-validation folds, we compute per-fold accuracy for models A and B, then test:

\begin{equation}
    H_0: \text{median}(d_i) = 0, \quad d_i = \text{acc}_A^{(i)} - \text{acc}_B^{(i)}
\end{equation}

Reject $H_0$ if $p < 0.05$.

\section{Experimental Setup}
\subsection{Hardware \& Software}
NVIDIA RTX 4090 (24 GB), Python 3.12, PyTorch 2.3, Ultralytics YOLO 8.2, FiftyOne 0.24, scikit-learn 1.4.

\subsection{Datasets}
Five YOLO-format datasets: industrial defect detection (250k images), satellite imagery (80k), medical cell counting (45k), autonomous driving (120k), agriculture pest detection (35k).

\subsection{Models}
YOLOv8n (3.2M params) vs. YOLO26n (2.8M params). Both trained for 250 epochs, imgsz=640, batch=16, lr0=0.01, 5-fold cross-validation.

\section{Results \& Discussion}

\subsection{Bootstrap Analysis}
\Cref{tab:bootstrap} shows the bootstrap CI for each dataset.

\begin{table}[htbp]
\centering
\caption{Bootstrap 95\% CI for mAP$_{50}$ Difference (YOLO26n $-$ YOLOv8n)}
\label{tab:bootstrap}
\begin{tabular}{@{}lcccc@{}}
\toprule
\textbf{Dataset} & \textbf{mAP Diff (\%)} & \textbf{CI Lower} & \textbf{CI Upper} & \textbf{Significant?} \\
\midrule
Defect Det. & +0.8 & +0.2 & +1.4 & \textbf{Yes} \\
Satellite & +0.3 & $-$0.4 & +1.0 & No \\
Medical Cells & +1.1 & +0.5 & +1.7 & \textbf{Yes} \\
Driving & +0.5 & $-$0.3 & +1.3 & No \\
Agriculture & +0.4 & $-$0.5 & +1.2 & No \\
\bottomrule
\end{tabular}
\end{table}

Of 5 apparent improvements, only 2 have CIs that exclude zero. The satellite, driving, and agriculture datasets show improvements within the noise floor.

\subsection{Failure Mode Analysis}
Across all datasets: 78\% of critical FPs involve class confusion (similar-looking defects), 15\% are localization errors, and 7\% are background confusion. Critical FPs average confidence 0.83---the model is confidently wrong.

\subsection{Wilcoxon Test}
$p$-values: Defect ($p = 0.008$), Medical ($p = 0.003$), Satellite ($p = 0.12$), Driving ($p = 0.21$), Agriculture ($p = 0.18$). Only Defect and Medical pass $\alpha = 0.05$.

\subsection{Ablation Study}
\begin{table}[htbp]
\centering
\caption{Impact of Statistical Validation on Claims}
\label{tab:ablation}
\begin{tabular}{@{}lc@{}}
\toprule
\textbf{Configuration} & \textbf{Overconfident Claims} \\
\midrule
No validation (baseline) & 3/5 (60\%) \\
Bootstrap only & 1/5 (20\%) \\
Wilcoxon only & 1/5 (20\%) \\
Full pipeline (Bootstrap + Wilcoxon) & 0/5 (0\%) \\
\bottomrule
\end{tabular}
\end{table}

\section{Data \& Code Availability Statement}
This architecture operates under a Dual Licensing Model (PolyForm Noncommercial / AGPLv3). To deploy the project and reproduce these experiments, use the \url{https://github.com/wisrovi/wyoloservice2\_production} repository.

\section{Broader Impact / Ethics Statement}
Overclaiming model improvements wastes researcher time and misleads practitioners. A 0.3\% mAP improvement that passes no statistical test should not be published as a contribution. Our pipeline makes statistical rigor automatic, reducing the barrier to responsible reporting.

\section{Conclusion \& Future Work}
We presented a statistical validation pipeline that combines bootstrap CIs, failure mode analysis, and Wilcoxon tests to determine whether YOLO model improvements are real. Of 5 apparent improvements, only 2 survived statistical scrutiny. Future work will extend the pipeline to support Bayesian model comparison and effect size estimation (Cohen's $d$).

\section{Acknowledgments}
We thank the contributors of the wisrovi-suit project for the foundational CLI and orchestration infrastructure.

\bibliographystyle{plain}
\bibliography{references}

\end{document}
